\documentclass[12pt]{article}
\usepackage{algorithmic}
\usepackage{algorithm}
\usepackage{amsbsy}
\usepackage{amsfonts}
\usepackage{amsmath}
\usepackage{amssymb}
\usepackage{amsthm}
\usepackage{array}
\usepackage{booktabs}
\usepackage{epsfig}
\usepackage{fancyhdr}
\usepackage{float}
\usepackage{fullpage}
\usepackage{graphicx}
\usepackage{ifthen}
\usepackage{times}
\usepackage{xspace}
\usepackage[colorlinks]{hyperref}
\usepackage[numbers,sort&compress]{natbib}
\usepackage[T1]{fontenc}
\usepackage{lmodern}

\usepackage{boldtensors}
\usepackage{microtype}

% Allow line breaks inside \texttt{...} at underscores and after "::"
% so long code identifiers (e.g. CapModel_Def.hpp, PHAL::AlbanyTraits::Jacobian)
% do not overrun the right margin.
\setlength{\emergencystretch}{3em}
\makeatletter
\def\BreakableUnderscore{\leavevmode\nobreak\hskip\z@{\char`\_}\nobreak\hskip\z@\allowbreak}
\makeatother
\AtBeginDocument{\def\_{\BreakableUnderscore}}
\newcommand{\origcolcol}{::}
\newcommand{\bcc}{::\allowbreak}

\DeclareMathOperator{\trace}{tr}
\DeclareMathOperator{\dev}{dev}
\DeclareMathOperator{\Grad}{Grad}
\DeclareMathOperator{\sym}{sym}
\newcommand{\transpose}[1]{{#1}^{\mathrm{T}}}
\newcommand{\Reals}{\mathbb{R}}

% Model-specific symbols
\newcommand{\stress}{~\sigma}
\newcommand{\kirch}{~\tau}
\newcommand{\backstress}{~\alpha}
\newcommand{\relstress}{~\xi}
\newcommand{\estrain}{{~\varepsilon^e}}
\newcommand{\pstrain}{{~\varepsilon^p}}
\newcommand{\strainincr}{\Delta~\varepsilon}
\newcommand{\Fdef}{~F}
\newcommand{\Fplas}{{~F^p}}
\newcommand{\bee}{{~b^e}}
\newcommand{\Celast}{{~{\mathsf C}^e}}
\newcommand{\ident}{~I}
\newcommand{\sdev}{~s}
\newcommand{\halpha}{~h^{\alpha}}

\begin{document}

\setlength{\headheight}{15pt}
\headsep = 4pt
\pagestyle{fancyplain}

\lhead
[\fancyplain{}{\emph{LCM Documentation}}]
{\fancyplain{}{\emph{LCM Documentation}}}

\rhead
[\fancyplain{}{\emph{Cap Plasticity}}]
{\fancyplain{}{\emph{Cap Plasticity}}}

\title{
  Cap Plasticity in LCM:\\
  Reference Formulation, Integration, and Verification}

\author{
  Alejandro Mota$^1$\thanks{Email: amota@sandia.gov}\\
  \\
  $^1$Mechanics of Materials Department\\
  Sandia National Laboratories\\
  Livermore CA 94550, USA\\
}

\date{\today}

\maketitle

\begin{abstract}

  LCM implements a single cap-plasticity constitutive model,
  \texttt{CapModel} (material-database name \texttt{Cap}): the
  three-invariant isotropic/kinematic hardening cap model of the
  Sandia GeoModel lineage
  \citep{Fossum.Fredrich:2000,Fossum.Brannon:2004,Foster.etal:2005},
  in the form presented in Section~3 of \citet{Sun.Chen.Ostien:2014}.
  This document states the formulation as implemented, including the
  conventions that the published papers leave ambiguous and that were
  settled against the original sources and the LAME GeoModel reference
  implementation; documents the adaptive substepped explicit
  integration; describes the logarithmic-strain / Kirchhoff-stress
  kinematics used for finite deformation; summarizes why no
  algorithmic tangent is ever hand-coded (Sacado automatic
  differentiation produces all of them); and describes the
  verification of the implementation against an independent reference
  implementation, which reproduces the kernel to machine precision
  over full load histories.
  \\
  \\
  Keywords: Cap plasticity, three-invariant yield surface, kinematic
  hardening, pore collapse, adaptive substepping, logarithmic strain,
  Kirchhoff stress, automatic differentiation, verification, LCM.

\end{abstract}

\section{Scope, Sources, and Conventions}
\label{sec:scope}

\texttt{CapModel} lives in
\texttt{src/LCM/models/CapModel.\{hpp,cpp\}} and
\texttt{CapModel\_Def.hpp} and is selected in the material database
with \texttt{Model Name: Cap}. It supports small-strain and
finite-deformation kinematics through the boolean material parameter
\texttt{Finite Deformation} (Section~\ref{sec:finite}); the
constitutive integrator is identical in both cases.

A fully implicit variant (\texttt{CapImplicit}, a 13-unknown local
backward-Euler solve following \citealt{Foster.etal:2005}) existed in
LCM until 2026 and was removed: it was unmaintained, had no test
coverage, and its principal classical advantage --- the consistent
algorithmic tangent --- carries little weight in LCM, where automatic
differentiation produces the exact tangent of whatever integrator
actually runs (Section~\ref{sec:ad}). The accuracy concern that
remains for an explicit scheme, large strain increments, is addressed
by adaptive substepping (Section~\ref{sec:integration}). The model's
own authors made the same choice, moving from the implicit
formulation of \citet{Foster.etal:2005} to the refined explicit
scheme of \citet{Sun.Chen.Ostien:2014} for production work.

\paragraph{Sources and their precedence.}
The implementation follows
\citet[Section~3]{Sun.Chen.Ostien:2014}. Where that paper is
ambiguous, in error, or silent, the conventions are settled by, in
order: the original model papers
\citep{Foster.etal:2005,Regueiro.Foster:2011}, and the LAME GeoModel
reference implementation (\texttt{iso\_geomodel\_model.F}, the
production Sandia code documented in \citealt{Fossum.Brannon:2004}).
Each such arbitration is called out where it appears below. Two
errata in \citet{Sun.Chen.Ostien:2014} as printed are relevant:
\begin{enumerate}
\item Algorithm~1, Step~2 prints the stress update as
      $\stress_{n+1} = \stress^{\mathrm{tr}}
       - \Delta\gamma\,\partial g/\partial\stress$, omitting the
      elastic tangent. The $\Delta\gamma$ derivation in the same
      paper, and Algorithm~2, both require
      $\stress_{n+1} = \stress^{\mathrm{tr}}
       - \Delta\gamma\,\Celast\!:\!\partial g/\partial\stress$,
      which is what LCM implements.
\item Algorithm~2's normal-correction fallback mixes $\Celast$ into
      the stress update while omitting it from the denominator
      $\tilde\chi = \partial_\sigma f : \partial_\sigma f$, which is
      dimensionally inconsistent. LCM implements the consistent
      \citet{Sloan.Abbo.Sheng:2001} form: a pure geometric projection
      with no elastic tangent in either place
      (Section~\ref{sec:integration}).
\end{enumerate}

\paragraph{Sign convention.}
Tension positive, compression negative, throughout. Under
confinement $I_1 = \trace\stress < 0$; the cap sits on the
compressive (negative-$I_1$) side, and ``the cap expands'' means its
hydrostat intercept moves toward $-\infty$.

\section{Physical Overview}
\label{sec:physics}

Geomaterials exhibit two inelastic mechanisms that the yield surface
must represent simultaneously. At low to moderate pressure the
dominant mechanism is frictional sliding and growth of microcracks:
shear resistance grows with confinement, producing a pressure-dependent
\emph{shear failure envelope} that, unlike a Drucker--Prager cone,
saturates once the cracks are squeezed shut. At high pressure a
porous material yields under hydrostatic load alone by irreversible
\emph{pore collapse}; representing this requires closing the yield
surface on the compressive end with a \emph{cap}, which moves outward
as porosity is crushed away --- the isotropic hardening of the model.
Shear stress assists pore collapse, so the cap is an ellipse rather
than a flat cutoff. Directionality of frictional sliding produces a
strong Bauschinger effect, represented by a deviatoric back stress
that translates the yield surface up to, but never beyond, the fixed
failure envelope. Finally, geomaterials are weaker in triaxial
extension than in triaxial compression, which brings the third stress
invariant into the yield function. The classical two-surface cap
models \citep[e.g.][]{Fossum.Fredrich:2000} join envelope and cap at
a corner; the present formulation is a single $C^1$ surface, which is
the property that makes the explicit integration well behaved.

\section{Kinematics in the Small-Strain Path}
\label{sec:smallstrain}

The small-strain path additively decomposes the strain rate
\citep[Eq.~47]{Sun.Chen.Ostien:2014},
\begin{equation}
  \dot{~\varepsilon} \;=\; \dot{\estrain} + \dot{\pstrain},
\end{equation}
and uses isotropic linear elasticity
\citep[Eqs.~49--50]{Sun.Chen.Ostien:2014},
\begin{equation}
  \dot{\stress} \;=\; \Celast : \dot{\estrain}
    \;=\; \Celast : (\dot{~\varepsilon} - \dot{\pstrain}),
  \qquad
  \Celast \;=\; \lambda~I \otimes ~I + 2\mu ~{\mathbb I},
\end{equation}
with $\lambda,\mu$ the Lam\'e constants and $~{\mathbb I}$ the
symmetric fourth-order identity. Given state
$\{\stress_n, \backstress_n, \kappa_n\}$ and strain increment
$\strainincr$ from the surrounding global Newton iteration, the trial
elastic stress is
\begin{equation}
  \stress^{\mathrm{tr}} \;=\; \stress_n + \Celast : \strainincr .
  \label{eq:trial_small}
\end{equation}

\section{Invariants, Lode Coefficient, and Yield Surface}
\label{sec:yield}

Kinematic hardening is introduced through the deviatoric back stress
$\backstress$ via the relative stress
\begin{equation}
  \relstress \;:=\; \stress - \backstress,
\end{equation}
with invariants
\begin{align}
  I_1        &= \trace \relstress, \\
  J_2^{\xi}  &= \tfrac{1}{2}\, \sdev : \sdev,
                \qquad \sdev \;=\; \relstress - \tfrac{1}{3} I_1 ~I, \\
  J_3^{\xi}  &= \det \sdev .
\end{align}
Because $\backstress$ is deviatoric,
$\trace\relstress = \trace\stress$, so the $\xi$ superscript is used
on $J_2, J_3$ only.

\paragraph{Lode coefficient.}
The difference in strength between triaxial compression and triaxial
extension enters through
\citep[Eqs.~34--36]{Sun.Chen.Ostien:2014}
\begin{equation}
  \Gamma(\beta) \;=\;
    \tfrac{1}{2}\left[
      1 - \sin 3\beta
    + \frac{1}{\psi}\left(1 + \sin 3\beta\right)
    \right],
  \qquad
  \sin 3\beta \;=\; -\frac{3\sqrt{3}\,J_3^\xi}{2\,(J_2^\xi)^{3/2}},
  \label{eq:Gamma}
\end{equation}
where $\psi \le 1$ is the ratio of triaxial-extension to
triaxial-compression strength. In triaxial compression $\Gamma = 1$;
in triaxial extension $\Gamma = 1/\psi$, and since $\Gamma^2$
multiplies $J_2^\xi$ in the yield function, yield in extension occurs
at $\psi$ times the compressive strength. The code sets
$\Gamma \equiv 1$ when $\psi = 0$ or $J_2^\xi = 0$ to avoid the
singularity on the hydrostat. With $\psi = 1$ the Lode dependence is
inactive and all $J_3$ terms drop out of the stress gradients ---
a fact relevant to the implementation history
(Appendix~\ref{sec:history}).

\paragraph{Shear failure surface, cap, and cap position.}
The pressure-dependent shear failure surface and the smooth
elliptical hardening cap
\citep[Eqs.~32, 37--38]{Sun.Chen.Ostien:2014} are
\begin{align}
  F_f(I_1) &= A - C \exp(D\,I_1) - \theta\,I_1,
             \label{eq:Ff}\\
  F_c(I_1,\kappa) &=
    1 - H(\kappa - I_1)\,\frac{(I_1 - \kappa)^2}{[X(\kappa) - \kappa]^2},
             \label{eq:Fc}\\
  X(\kappa) &= \kappa - R\,F_f(\kappa),
             \label{eq:X}
\end{align}
where $H$ is the Heaviside function. The internal variable $\kappa$
is the \emph{branch point} where the cap begins, and $X(\kappa)$ is
the \emph{hydrostat intercept}: the current hydrostatic crush
strength. On the shear branch $I_1 > \kappa$, $F_c \equiv 1$; on the
compactive cap $X < I_1 \le \kappa$, $F_c$ falls quadratically from
one to zero, closing the surface at $I_1 = X$. At $I_1 = \kappa$
both $F_c = 1$ and $\partial F_c/\partial I_1 = 0$, so the surface is
$C^1$ across the branch point: there is no corner where cap meets
envelope. The strength at zero pressure is $A - C$; as compression
grows the exponential dies and strength saturates toward
$A - \theta I_1$. $R$ is the cap aspect ratio (its extent in $I_1$
per unit of $F_f$), and $N$ offsets the initial yield surface below
the failure envelope, leaving room for kinematic hardening.

\paragraph{Yield function and plastic potential.}
\begin{align}
  f(I_1, J_2^\xi, J_3^\xi, \backstress, \kappa)
    &= \Gamma^2 J_2^\xi - F_c(I_1,\kappa)\,\bigl(F_f(I_1) - N\bigr)^2,
    \label{eq:f}\\
  g(I_1, J_2^\xi, J_3^\xi, \backstress, \kappa)
    &= \Gamma^2 J_2^\xi - F_c^g(I_1,\kappa)\,\bigl(F_f^g(I_1) - N\bigr)^2,
    \label{eq:g}
\end{align}
with the plastic-potential analogues
\begin{align}
  F_f^g(I_1)  &= A - C \exp(L\,I_1) - \phi\,I_1, \label{eq:Ffg}\\
  F_c^g(I_1,\kappa) &=
    1 - H(\kappa - I_1)\,\frac{(I_1 - \kappa)^2}{[X^g(\kappa) - \kappa]^2},\\
  X^g(\kappa) &= \kappa - Q\,F_f(\kappa). \label{eq:Xg}
\end{align}
Non-associativity ($L \ne D$, $\phi \ne \theta$, $Q \ne R$) exists
because associated flow on a frictional envelope grossly
over-predicts dilatancy; flattening the potential's pressure
dependence ($\phi < \theta$, $L < D$) reduces the predicted plastic
volume expansion to measured levels. Associative flow corresponds to
$L = D$, $\phi = \theta$, $Q = R$.

\emph{Convention note.} In Eq.~\eqref{eq:Xg} the potential cap
position $X^g$ is built from the \emph{yield} failure function $F_f$
(parameters $D, \theta$), not from $F_f^g$: only the aspect ratio
$Q$ distinguishes the potential cap from the yield cap. This follows
\citet[Eq.~14]{Regueiro.Foster:2011} and the LAME reference
implementation. The distinction is invisible for associative
parameters and matters otherwise; the pre-2026 LCM code had it wrong
(Appendix~\ref{sec:history}).

\section{Flow Rule and Hardening Evolution}
\label{sec:flow}

The plastic strain rate follows the gradient of the plastic
potential \citep[Eq.~48]{Sun.Chen.Ostien:2014},
\begin{equation}
  \dot{\pstrain} \;=\; \dot\gamma\,\frac{\partial g}{\partial \stress},
  \label{eq:flow}
\end{equation}
where $\dot\gamma \ge 0$ is the consistency parameter.

\paragraph{Kinematic hardening.}
The deviatoric back stress evolves along the deviator of the flow
direction \citep[Eqs.~41--43]{Sun.Chen.Ostien:2014},
\begin{equation}
  \dot{\backstress} \;=\; \dot\gamma\,\halpha,
  \qquad
  \halpha \;=\; c^{\alpha}\,G^{\alpha}(\backstress)\,
               \dev\!\left(\frac{\partial g}{\partial \stress}\right),
  \qquad
  G^{\alpha} \;=\;
    \max\!\left(0,\; 1 - \frac{\sqrt{J_2^{\alpha}}}{N}\right),
  \label{eq:halpha}
\end{equation}
with $J_2^{\alpha} = \tfrac{1}{2}\backstress:\backstress$. The
limiter $G^\alpha$ drives the hardening rate to zero as the
translating yield surface approaches the fixed failure envelope,
which it reaches when $\sqrt{J_2^\alpha} = N$.

\emph{The clamp at zero in Eq.~\eqref{eq:halpha} is essential and
absent from the published papers.} The explicit update can overshoot
$\sqrt{J_2^\alpha}$ slightly past $N$; without the clamp,
$G^\alpha < 0$ \emph{reverses} the direction of kinematic hardening.
At a material point this is a small error, but at the structural
level it destabilizes the homogeneous solution: in LCM's verification
study a uniform-property confined-compression problem on a
$2 \times 2 \times 2$ mesh demonstrably bifurcated to a
non-homogeneous Newton solution at exactly the step where the back
stress reached its limit surface. The LAME reference implementation
has always floored its equivalent quantity (\texttt{GFUN}) at zero.
See Appendix~\ref{sec:history}.

\paragraph{Isotropic (cap) hardening.}
The cap position is driven by volumetric plastic strain through the
experimentally measured crush curve
\citep[Eqs.~44--46]{Sun.Chen.Ostien:2014},
\begin{equation}
  \dot\kappa \;=\; \dot\gamma\,h^{\kappa},
  \qquad
  h^{\kappa} \;=\;
    \frac{\trace(\partial g/\partial \stress)}
         {\partial \varepsilon_v^p/\partial \kappa},
  \qquad
  \varepsilon_v^p(\kappa)
    = W \Bigl[
        e^{\,D_1 \bar X - D_2 \bar X^2} - 1
      \Bigr],
  \quad \bar X = X(\kappa) - X(\kappa_0),
  \label{eq:kappa}
\end{equation}
where $\kappa_0$ is the initial cap branch point and
$W, D_1, D_2$ are material parameters. $W$ is the maximum
compactive volumetric plastic strain, essentially the initial
porosity: as $X \to -\infty$, $\varepsilon_v^p \to -W$ and all
porosity has been crushed out. The denominator in $h^\kappa$ is
evaluated by the chain rule
$(\partial\varepsilon_v^p/\partial X)(\mathrm{d}X/\mathrm{d}\kappa)$
in \texttt{compute\_dedkappa}.

\emph{Convention note.} The crush curve in Eq.~\eqref{eq:kappa} is
parametrized by the \emph{yield} cap position $X(\kappa)$ of
Eq.~\eqref{eq:X} --- parameters $R, D, \theta$ --- because $X$ is the
measurable hydrostatic elastic limit against which the curve is
calibrated. The LAME reference computes
$\mathrm{d}X/\mathrm{d}\kappa = 1 - R\,\mathrm{d}F_f/\mathrm{d}\kappa$
accordingly. The pre-2026 LCM code used the potential-surface
quantities ($Q, L, \phi$) here, which is wrong for non-associative
parameters (Appendix~\ref{sec:history}).

Two safeguards from the LAME reference, absent from the papers,
apply wherever $\kappa$ is updated:
\begin{itemize}
\item \emph{No cap contraction:} increments of $\kappa$ are clamped
      to $\Delta\kappa \le 0$ (LAME:
      \texttt{HK = MIN(HK,0)}, comment ``prevents cap
      contraction''). Dilation never re-opens porosity.
\item \emph{Cap lock at full compaction:} when
      $|\varepsilon_v^p(\kappa)| \ge W$ the crush curve is exhausted,
      its slope vanishes, and $h^\kappa$ as written diverges.
      Following the reference, $h^\kappa$ is then set to
      $-0.01\,K^2$ ($K$ the bulk modulus), which races the cap toward
      $-\infty$ so that pore-collapse plasticity shuts off rather
      than stalling against an immobile cap.
\end{itemize}

\section{Integration: Substepped Explicit Scheme}
\label{sec:integration}

The integrator is the refined explicit scheme of
\citet[Algorithms~1--2]{Sun.Chen.Ostien:2014}, upgraded with the
adaptive substepping of \citet{Sloan.Abbo.Sheng:2001}: within the
plastic part of a strain increment, modified-Euler (RK1/RK2) pairs
advance the state, the difference between the two estimates supplies
a local error measure that controls substep size, and every accepted
substep is followed by a drift correction back to the yield surface.

\paragraph{Rate evaluation.}
At state $\{\stress, \backstress, \kappa\}$, for a substrain
$\strainincr_{\mathrm{sub}}$, the consistency condition
$\dot f = 0$ gives
\begin{equation}
  \Delta\gamma \;=\;
  \frac{(\partial f/\partial \stress) : \Celast :
        \strainincr_{\mathrm{sub}}}{\chi},
  \qquad
  \chi \;=\; \tfrac{\partial f}{\partial \stress}
            :\Celast:\tfrac{\partial g}{\partial \stress}
          - \tfrac{\partial f}{\partial \backstress}:\halpha
          - \tfrac{\partial f}{\partial \kappa}\,h^{\kappa},
  \label{eq:dgamma}
\end{equation}
with $\partial f/\partial\backstress = -\partial f/\partial\stress$.
$\Delta\gamma$ is clamped at zero (elastic unloading within a
substep), and the increments are
\begin{equation}
  \Delta\stress = \Celast : \strainincr_{\mathrm{sub}}
                - \Delta\gamma\,\Celast : \tfrac{\partial g}{\partial\stress},
  \qquad
  \Delta\backstress = \Delta\gamma\,\halpha,
  \qquad
  \Delta\kappa = \min(0,\,\Delta\gamma\,h^{\kappa}).
  \label{eq:rates}
\end{equation}

\paragraph{Algorithm.}
With input state $\{\stress_n, \backstress_n, \kappa_n\}$ and strain
increment $\strainincr$:
\begin{enumerate}
\item \emph{Elastic trial} via~\eqref{eq:trial_small}. If
      $f(\stress^{\mathrm{tr}}, \backstress_n, \kappa_n) \le 0$,
      accept the trial and stop.
\item \emph{Substepping}, starting from the converged state (not the
      trial), with pseudo-time $T \in [0,1]$, $\Delta T = 1$
      initially, and $\Delta T_{\min} = 1/\texttt{Maximum Substeps}$.
      While $T < 1$:
      \begin{enumerate}
      \item Evaluate the increments \eqref{eq:rates} at the current
            state for $\strainincr_{\mathrm{sub}}
            = \Delta T\,\strainincr$ (RK1), and again at the
            RK1-advanced state (RK2).
      \item Form the modified-Euler state from the average of the two
            increment sets, and the relative error estimate
            \begin{equation}
              R \;=\; \frac{\|\Delta\stress_2 - \Delta\stress_1\|}
                           {2\,\max(\|\stress_{\mathrm{new}}\|,\,
                            \epsilon)} .
            \end{equation}
      \item If $R > \texttt{Substep Tolerance}$ and
            $\Delta T > \Delta T_{\min}$: reject, set
            $\Delta T \leftarrow
            \max\bigl(0.9\sqrt{\mathrm{STOL}/R},\,0.1\bigr)\,\Delta T$
            (floored at $\Delta T_{\min}$), retry.
      \item Otherwise accept: apply the drift correction below,
            advance $T \leftarrow T + \Delta T$, and grow
            $\Delta T \leftarrow
            \min\bigl(0.9\sqrt{\mathrm{STOL}/R},\,1.1\bigr)\,\Delta T$,
            capped at $1 - T$.
      \end{enumerate}
\item \emph{Plastic strain increment} recovered through elastic
      compliance,
      \begin{equation}
        \Delta\pstrain \;=\; (\Celast)^{-1} : (\stress^{\mathrm{tr}}
                                              - \stress_{n+1}),
      \end{equation}
      from which the outputs
      $\varepsilon_v^p \mathrel{+}= \trace\Delta\pstrain$ and
      $\varepsilon_{\mathrm{eqps}} \mathrel{+}=
      \sqrt{\tfrac{2}{3}}\,\|\dev\Delta\pstrain\|$ are accumulated.
\end{enumerate}

\paragraph{Drift correction.}
Explicit updates do not satisfy $f = 0$ exactly. After each accepted
substep, while $|f| > f_{\mathrm{tol}}$ and fewer than 20 iterations:
\begin{enumerate}
\item \emph{Consistent correction:} $\delta\gamma = f/\chi$ with all
      quantities at the current state; update
      $\stress \leftarrow \stress
       - \delta\gamma\,\Celast:\partial g/\partial\stress$,
      $\backstress \leftarrow \backstress + \delta\gamma\,\halpha$,
      $\kappa \leftarrow \kappa + \min(0, \delta\gamma\,h^\kappa)$.
\item \emph{Normal-correction fallback:} if the consistent correction
      increases $|f|$, discard it and instead project geometrically,
      \begin{equation}
        \delta\gamma \;=\;
          \frac{f}{\partial_\sigma f : \partial_\sigma f},
        \qquad
        \stress \leftarrow \stress
          - \delta\gamma\,\frac{\partial f}{\partial\stress},
      \end{equation}
      holding $\backstress, \kappa$ fixed. This is the
      \citet{Sloan.Abbo.Sheng:2001} form; the version printed in
      \citet{Sun.Chen.Ostien:2014} is dimensionally inconsistent
      (Section~\ref{sec:scope}).
\end{enumerate}
The yield function has units of stress squared, so the drift
tolerance is scaled to make the convergence test independent of the
unit system:
\begin{equation}
  f_{\mathrm{tol}} \;=\; 10^{-12}\,E^2 .
\end{equation}
(The historical absolute tolerance of $10^{-10}$ was unreachable in
Pa-based unit systems, so the corrector always ran to its iteration
cap.)

\paragraph{Branch consistency under AD.}
The substep control flow (acceptance, rejection, step growth)
branches on Sacado \emph{values}, so the Residual and Jacobian
evaluations of the same state take identical branch sequences and the
AD tangent differentiates exactly the code path that produced the
stress.

\paragraph{Material parameters of the integrator.}
\texttt{Substep Tolerance} (STOL, default $10^{-4}$) and
\texttt{Maximum Substeps} (default 200). A rejected attempt counts
toward an overall attempt cap of $4\times$\texttt{Maximum Substeps}.

\paragraph{Hand-coded derivatives.}
The closed-form constitutive derivatives
$\partial f/\partial \stress$, $\partial g/\partial \stress$,
$\partial f/\partial \kappa$, $\halpha$, and
$\partial\varepsilon_v^p/\partial\kappa$ are provided in
\texttt{CapModel\_Def.hpp} (\texttt{compute\_dfdsigma},
\texttt{compute\_dgdsigma}, \texttt{compute\_dfdkappa},
\texttt{compute\_halpha}, \texttt{compute\_dedkappa}). The only
nontrivial tensor derivative among them is
\begin{equation}
  \frac{\partial J_3}{\partial \stress}
  \;=\; \sdev\cdot\sdev - \tfrac{2}{3} J_2\,\ident ,
  \label{eq:dJ3}
\end{equation}
a matrix product, \emph{not} the elementwise square of $\sdev$; the
two coincide only in principal axes (Appendix~\ref{sec:history}).
Every one of these closed forms is checked against complex-step
differentiation by the verification suite
(Section~\ref{sec:verification}).

\section{Finite-Deformation Kinematics}
\label{sec:finite}

With \texttt{Finite Deformation: true} the model uses the
multiplicative split and the isotropic logarithmic-strain /
Kirchhoff-stress framework
\citep{Simo:1992}: because every constitutive function in
Sections~\ref{sec:yield}--\ref{sec:flow} is an isotropic function of
its tensor arguments, the small-strain integrator of
Section~\ref{sec:integration} runs \emph{unchanged} with the
Kirchhoff stress $\kirch$ and the logarithmic elastic strain in the
slots that otherwise hold $\stress$ and $~\varepsilon$. Only the
kinematic pre- and post-processing differ. At each quadrature point,
with incoming $\Fdef$, stored $\Fplas_n$, and the deformation
gradient $\Fdef_n$ of the previous converged step:
\begin{enumerate}
\item \emph{Trial elastic left Cauchy--Green tensor} from the stored
      plastic configuration
      ($~C_p^{-1} = (\Fplas_n)^{-1}(\Fplas_n)^{-\mathrm{T}}$),
      \begin{equation}
        \bee^{\mathrm{tr}} \;=\; \Fdef\,~C_p^{-1}\,\transpose{\Fdef},
        \qquad
        \estrain^{\mathrm{tr}} \;=\;
        \tfrac{1}{2}\,\log_{\mathrm{sym}}\bee^{\mathrm{tr}},
      \end{equation}
      with $\log_{\mathrm{sym}}$ the principal (spectral) logarithm.
\item \emph{Starting stress and strain increment.} The elastic
      logarithmic strain at $t_n$ is recovered from the same plastic
      configuration and the old deformation gradient,
      \begin{equation}
        \estrain_n \;=\;
        \tfrac{1}{2}\,\log_{\mathrm{sym}}
        \bigl(\Fdef_n\,~C_p^{-1}\,\transpose{(\Fdef_n)}\bigr),
        \qquad
        \kirch_n = \Celast : \estrain_n,
        \qquad
        \strainincr = \estrain^{\mathrm{tr}} - \estrain_n .
      \end{equation}
      The integrator then starts from $\kirch_n$ with increment
      $\strainincr$, so its trial stress
      $\kirch_n + \Celast:\strainincr = \Celast:\estrain^{\mathrm{tr}}$
      is \emph{exact} for the hyperelastic logarithmic model: the
      incremental form introduces no approximation in the elastic
      relation.
\item \emph{Return mapping} of Section~\ref{sec:integration} on
      $\kirch$, yielding $\kirch_{n+1}, \backstress_{n+1},
      \kappa_{n+1}$.
\item \emph{Plastic configuration update.} The new elastic
      logarithmic strain follows from the returned stress,
      $\estrain_{n+1} = (\Celast)^{-1}:\kirch_{n+1}$, whence
      \begin{equation}
        \bee_{n+1} = \exp_{\mathrm{sym}}(2\,\estrain_{n+1}),
        \qquad
        ~C_{p,n+1}^{-1} = \Fdef^{-1}\,\bee_{n+1}\,\Fdef^{-\mathrm{T}},
        \qquad
        \Fplas_{n+1} = \bigl(~C_{p,n+1}\bigr)^{1/2},
      \end{equation}
      where the last expression is the symmetric positive-definite
      square root. The rotational part of $\Fplas$ is unobservable in
      an isotropic model, so the symmetric root is a legitimate
      representative; no exponential map of the flow direction is
      formed explicitly --- it is absorbed in the spectral
      $\log$/$\exp$ of $\bee$.
\item \emph{Cauchy stress} for assembly,
      $\stress_{n+1} = \kirch_{n+1}/J$ with $J = \det\Fdef$.
\end{enumerate}
The stored state in this mode is $\Fplas$ (identity-initialized)
plus the deformation gradient history that
\texttt{MechanicsProblem} already maintains; the small-strain mode
stores the strain tensor instead. All other state (stress, back
stress, $\kappa$, $\varepsilon_{\mathrm{eqps}}$, $\varepsilon_v^p$)
is common.

This wrapper is independent of how the inner return mapping is
integrated, which is why the desire for finite-deformation support
had no bearing on the decision to remove the implicit integrator.

\section{Tangents Are Not Hand-Coded: Automatic Differentiation}
\label{sec:ad}

\citet{Sun.Chen.Ostien:2014} do not document a consistent algorithmic
tangent $\partial \stress_{n+1}/\partial ~\varepsilon_{n+1}$, nor a
finite-deformation counterpart. LCM does not need either. LCM's
evaluators are templated on an evaluation trait; for
\texttt{PHAL\bcc AlbanyTraits\bcc Jacobian} every field entering the
kernel --- strain or deformation gradient, stored state, every
intermediate of the return mapping --- carries a Sacado forward-mode
AD type whose active directions are the element-level global degrees
of freedom. The entire integrator (trial, substepping, drift
correction, finite-deformation pre/post-processing) executes in AD
arithmetic, and the derivative of the output stress with respect to
every element DOF emerges as a by-product. Because the substep
control flow branches on Sacado values
(Section~\ref{sec:integration}), the differentiated path is exactly
the primal path. Sensitivities with respect to material parameters
(for calibration and inverse problems) are generated by the same
mechanism. The only hand-coded derivatives in the model are the
\emph{constitutive} ones listed in Section~\ref{sec:integration} ---
part of the model's definition, not of its integrator.

\section{Verification}
\label{sec:verification}

The implementation is verified, not merely regression-tested, by the
infrastructure in \texttt{tests/LCM/CapModelPlasticity3D}:

\paragraph{Independent reference implementation.}
\texttt{cap\_reference.py} is a deliberately separate transcription
of the model from the publications --- not a port of the C++. It
evaluates all stress gradients by complex-step differentiation
(machine-exact for the analytic branches), so the C++ closed-form
derivatives and the reference share no code and no derivation. Its
self-check verifies, among other identities, the $J_3$ derivative
\eqref{eq:dJ3} against complex step at a generic (non-principal-axes)
state, where the elementwise-square error would be $O(10^2)$ in
stress-squared units.

\paragraph{Material-point exactness of the test problems.}
Each verification problem is a single hexahedral element whose 24
degrees of freedom are all prescribed by face boundary conditions.
With no free degrees of freedom, the finite-element problem \emph{is}
the material point: no equilibrium iteration intervenes between the
prescribed strain history and the constitutive response, and the
comparison with the reference is exact rather than approximate. (The
importance of this is not academic: on a $2\times2\times2$ mesh the
same problem bifurcated away from the homogeneous solution when the
unclamped $G^\alpha$ defect of Appendix~\ref{sec:history} was
present.)

\paragraph{Load paths and results.}
Five paths are compared state-by-state (stress components, $\kappa$,
$\varepsilon_v^p$) over their full histories against the reference,
which runs the identical substepping algorithm:
\begin{center}
\small
\begin{tabular}{lp{3.4cm}p{4.4cm}l}
\toprule
Test & Path & Exercises & Max.\ rel.\ diff. \\
\midrule
\texttt{hydrostatic} & $~\varepsilon = -0.02\,t\,\ident$ &
  cap only; analytic identities & $\sim 10^{-14}$ \\
\texttt{confined} & $\varepsilon_{xx} = -0.04\,t$, 200 steps &
  shear--cap interaction, $\backstress$ to limit & $\sim 10^{-15}$ \\
\texttt{triaxial} & unequal normal strains &
  $\psi = 0.8$ (Lode), non-associative & $\sim 10^{-15}$ \\
\texttt{confined\_coarse} & confined, 20 steps &
  substep rejection/regrowth & $\sim 10^{-15}$ \\
\texttt{confined\_fd}, \texttt{triaxial\_fd} & as above, finite deformation &
  $\log$/$\exp$ kinematics & $\sim 10^{-12}$ \\
\bottomrule
\end{tabular}
\end{center}
The hydrostatic path additionally checks two analytic identities: on
the hydrostat the converged stress must satisfy $I_1 = X(\kappa)$
(observed to $\sim 10^{-13}$), and the accumulated
$\varepsilon_v^p$ must follow the crush curve of
Eq.~\eqref{eq:kappa} (observed to the per-step consistency error of
the discretization). The finite-deformation comparisons carry the
slightly larger $\sim 10^{-12}$ floor of the eigendecomposition
round-trips. The comparisons require \texttt{python3} with
\texttt{numpy} and \texttt{netCDF4} and are skipped with a status
message when these are unavailable.

\section{Material Parameters}
\label{sec:parameters}

All parameters are read from the material database with the names
below. The Salem limestone values are the associative set of
\citet[Table~1]{Sun.Chen.Ostien:2014}; stress-like quantities in MPa.

\begin{center}
\small
\begin{tabular}{llrp{6.4cm}}
\toprule
Input name & Symbol & Salem & Meaning \\
\midrule
\texttt{A}       & $A$        & 689.2    & saturated shear strength \\
\texttt{C}       & $C$        & 675.2    & strength deficit at zero pressure ($A{-}C$ = cohesion-like intercept) \\
\texttt{D}       & $D$        & 3.94e-4  & pressure scale of the exponential, yield surface \\
\texttt{theta}   & $\theta$   & 0.0      & linear pressure slope, yield surface \\
\texttt{N}       & $N$        & 6.0      & offset of initial yield below failure envelope \\
\texttt{R}       & $R$        & 28.0     & cap aspect ratio, yield surface \\
\texttt{kappa0}  & $\kappa_0$ & $-8.05$  & initial cap branch point \\
\texttt{W}       & $W$        & 0.08     & maximum compactive volumetric plastic strain ($\approx$ porosity) \\
\texttt{D1}      & $D_1$      & 1.47e-3  & crush-curve shape (linear) \\
\texttt{D2}      & $D_2$      & 0.0      & crush-curve shape (quadratic) \\
\texttt{calpha}  & $c^\alpha$ & 1.0e5    & kinematic hardening rate \\
\texttt{psi}     & $\psi$     & 1.0      & triaxial extension/compression strength ratio \\
\texttt{L}       & $L$        & 3.94e-4  & pressure scale, plastic potential \\
\texttt{phi}     & $\phi$     & 0.0      & linear pressure slope, plastic potential \\
\texttt{Q}       & $Q$        & 28.0     & cap aspect ratio, plastic potential \\
\midrule
\texttt{Substep Tolerance} & STOL & $10^{-4}$ &
  relative stress-error target per substep \\
\texttt{Maximum Substeps}  & ---  & 200 &
  smallest substep is $1/$this fraction of the increment \\
\texttt{Finite Deformation} & --- & false &
  $\log$/Kirchhoff kinematics of Section~\ref{sec:finite} \\
\bottomrule
\end{tabular}
\end{center}

For orientation with the Salem values: the initial hydrostatic crush
strength is $X_0 = \kappa_0 - R\,F_f(\kappa_0) \approx -460$~MPa,
i.e.\ pore collapse under pure pressure begins near
$p \approx 153$~MPa, while the cap branch point sits at only
$8$~MPa --- a long, flat cap, typical for porous rock.

\appendix

\section{Revision History of the Formulation}
\label{sec:history}

In 2026 the implementation was audited line-by-line against the
sources listed in Section~\ref{sec:scope}, treating the code as
unreliable, and then verified against the independent reference of
Section~\ref{sec:verification}. Results produced with the pre-audit
code differ from the current model in the following ways; all four
defects are absent from every published validation case, which is
why they survived.

\begin{enumerate}
\item \emph{$\partial J_3/\partial\stress$ used the elementwise
      square} $s_{ij} s_{ij}$ (no sum) instead of the matrix product
      $(\sdev\cdot\sdev)_{ij}$ in Eq.~\eqref{eq:dJ3}. The two agree
      in principal axes and the term is multiplied by
      $\mathrm{d}\Gamma/\mathrm{d}J_3 \propto (1 - 1/\psi)$, so any
      run with $\psi = 1$ --- including all of
      \citet{Sun.Chen.Ostien:2014} --- never exercised it. Runs with
      $\psi \ne 1$ received wrong plastic flow directions off the
      principal axes.
\item \emph{$X^g$ was built from $F_f^g$} ($L, \phi$) instead of the
      yield $F_f$ ($D, \theta$) required by
      \citet[Eq.~14]{Regueiro.Foster:2011}; see Eq.~\eqref{eq:Xg}.
      Invisible for associative parameters.
\item \emph{The crush-curve derivative used the potential-surface
      cap position} ($Q, L, \phi$) instead of the yield $X(\kappa)$
      ($R, D, \theta$); see Eq.~\eqref{eq:kappa}. Invisible for
      associative parameters.
\item \emph{$G^\alpha$ was unclamped}, allowing hardening reversal
      after overshoot of the back-stress limit surface; see
      Eq.~\eqref{eq:halpha}. This is the likely root cause of the
      Newton divergence that led to the disabling of the original
      \texttt{CapModelPlasticity3D} regression tests years earlier.
\end{enumerate}

In the same effort: the drift tolerance was made unit-system
invariant ($10^{-12} E^2$); the $\Delta\kappa \le 0$ clamp, present
in the code but undocumented, was confirmed against the LAME
reference and documented; the cap lock at full compaction was added
from the reference; adaptive substepping was added; the implicit
variant (\texttt{CapImplicit}) was removed; the model was renamed
from \texttt{CapExplicit} to \texttt{CapModel} with material-database
name \texttt{Cap}; and the earlier finite-deformation path (which
skipped the predictor, disabled the normal-correction fallback, and
updated $\Fplas$ by an exponential map of the final flow direction)
was replaced by the formulation of Section~\ref{sec:finite}.

\begin{thebibliography}{9}

\bibitem[Fossum and Brannon(2004)]{Fossum.Brannon:2004}
A.~F. Fossum and R.~M. Brannon.
\newblock {\em The Sandia GeoModel: Theory and User's Guide}.
\newblock Sandia Technical Report SAND2004-3226. Sandia National Laboratories,
  2004.

\bibitem[Fossum and Fredrich(2000)]{Fossum.Fredrich:2000}
A.~F. Fossum and J.~T. Fredrich.
\newblock {\em Cap plasticity models and compactive and dilatant pre-failure
  deformation}.
\newblock Sandia Technical Report SAND2000-1056. Sandia National Laboratories,
  2000.

\bibitem[Foster et~al.(2005)Foster, Regueiro, Fossum, and
  Borja]{Foster.etal:2005}
C.~D. Foster, R.~A. Regueiro, A.~F. Fossum, and R.~I. Borja.
\newblock Implicit numerical integration of a three-invariant,
  isotropic/kinematic hardening cap plasticity model for geomaterials.
\newblock {\em Computer Methods in Applied Mechanics and Engineering},
  194:5109--5138, 2005.

\bibitem[Regueiro and Foster(2011)]{Regueiro.Foster:2011}
R.~A. Regueiro and C.~D. Foster.
\newblock Bifurcation analysis for a rate-sensitive, non-associative,
  three-invariant, isotropic/kinematic hardening cap plasticity model for
  geomaterials: Part I. Small strain.
\newblock {\em International Journal for Numerical and Analytical Methods in
  Geomechanics}, 35:201--225, 2011.

\bibitem[Simo(1992)]{Simo:1992}
J.~C. Simo.
\newblock Algorithms for static and dynamic multiplicative plasticity that
  preserve the classical return mapping schemes of the infinitesimal theory.
\newblock {\em Computer Methods in Applied Mechanics and Engineering},
  99(1):61--112, 1992.

\bibitem[Sloan et~al.(2001)Sloan, Abbo, and Sheng]{Sloan.Abbo.Sheng:2001}
S.~W. Sloan, A.~J. Abbo, and D.~Sheng.
\newblock Refined explicit integration of elastoplastic models with automatic
  error control.
\newblock {\em Engineering Computations}, 18(1/2):121--154, 2001.

\bibitem[Sun et~al.(2014)Sun, Chen, and Ostien]{Sun.Chen.Ostien:2014}
W.~Sun, Q.~Chen, and J.~T. Ostien.
\newblock Modeling the hydro-mechanical responses of strip and circular punch
  loadings on water-saturated collapsible geomaterials.
\newblock {\em Acta Geotechnica}, 9(6):903--934, 2014.
\newblock doi:10.1007/s11440-013-0276-x.

\end{thebibliography}

\end{document}
